% p.62 — الباب الثامن والعشرون في كل الافعال المستعملة (تتمة)
\mspage{62}
\nohairlines
% Layout (as in the lithograph): one boxed table of seven narrow columns, read right-to-left,
% each column top-to-bottom; every cell is «رقم + فعل», with group headings (ملحق بالرباعي …,
% باب فتح …) as cells inside the columns.  Printed page number ۶۲ at top left.  The columns
% have unequal numbers of rows (19/21/18/21/21/21/7); empty cells below are padding.  The
% leftmost column holds, under its seven verbs, a short Arabic prose colophon, set here as a
% \multirow cell over the remaining 14 rows of that column.  Columns 1 and 7 are p-columns so
% that the long group headings and the colophon wrap; the whole table is scaled to the width.
% TODO: 61 — the lithograph writes the initial hamza/madda with the same stroke; read as آتٰى.
% TODO: 82 — the last letter of the verb is unclear (ح/ج/خ?).
\begin{sarf}{|>{\centering\arraybackslash}p{24mm}|r|r|r|r|r|>{\centering\arraybackslash}p{26mm}|}
(5) خَيْعَلَ & (28) قَالَ & (34) جَاءَ & \thead{باب حَسِبَ} & (61) \unclear{آتٰى} & (77) اِنْقَادَ & (96) اِسْتَوْجَبَ\\
(6) شَرْيَفَ & (35) دَعَا & (40) رَمٰى & (19) وَرِمَ & (62) أَرْأٰى & (78) اِنْطَارَ & (97) اِسْتَيْقَظَ\\
(7) قَلْسٰى & (38) كَنٰى & (42) أَرٰى & (25) يَقِظَ & (63) أَوْحٰى & (79) اِنْمَحٰى & (98) اِسْتَعَانَ\\
\thead{ملحق بالرباعي المزيد تَدَحْرَجَ} & \thead{باب فَتَحَ} & (45) وَقٰى & (47) وَلِيَ & (64) آوٰى & (80) اِنْقَضٰى & (99) اِسْتَمَالَ\\
(1) تَجَلْبَبَ & (2) أَهَبَ & (48) طَوٰى & \thead{باب كَرُمَ} & \thead{باب فَعَّلَ} & (81) اِنْزَوٰى & (100) اِسْتَحْلٰى\\
(2) تَقَلْنَسَ & (6) سَأَلَ & (50) أَوٰى & (5) أَدُبَ & (65) زَكّٰى & \thead{باب اِفْتَعَلَ} & (101) اِسْتَرْقٰى\\
(3) تَجَوْرَبَ & (9) قَرَأَ & (51) وَأٰى & (8) لَؤُمَ & (66) سَمّٰى & (82) \unclear{اِيْتَلَحَ} & (102) اِسْتَوْفٰى\\\cline{7-7}
(4) تَسَرْوَلَ & (16) وَقَعَ & \thead{باب سَمِعَ} & (12) جَرُؤَ & \thead{باب فَاعَلَ} & (83) اِهْتَمَّ & \multirow{14}{26mm}{هذه الافعال المستعملة المشهورة وَتَرَكْتُ الشُّوَاذَ الغير المعروفة، يمكن ان تُؤْخَذَ من هذه الافعال، يلزم على الاستاذ ان يمرنها حتى يَرْسُخَ في القلوب حتى يكون الاخذ سهلا جدا. تمَّ الصرف بعون الله ويليه النحو ان شاء الله الوهاب.}\\
(5) تَخَيْعَلَ & (22) يَنَعَ & (4) أَرِجَ & (20) وَسُمَ & (67) حَابَّ & (84) اِتَّقَدَ & \\
(6) تَشَرْيَفَ & (29) بَاءَ & (7) يَئِسَ & (26) يَمُنَ & (68) رَاعٰى & (85) اِتَّسَرَ & \\
(7) تَقَلْسٰى & (31) شَاءَ & (11) \unclear{سَبِأَ} & (37) رَخُوَ & (69) لَاقٰى & (86) اِجْتَابَ & \\
\thead{ملحق بالرباعي المزيد اِحْرَنْجَمَ} & (39) سَعٰى & (15) مَسَّ & \thead{باب أَفْعَلَ} & (70) آوٰى & (87) اِخْتَارَ & \\
(1) اِقْعَنْسَسَ & (44) رَأٰى & (18) وَجِلَ & (52) آمَنَ & \thead{باب تَفَعَّلَ} & (88) اِرْتَضٰى & \\
(2) اِحْوَنْصَلَ & \thead{باب ضَرَبَ} & (24) يَبِسَ & (53) أَحَبَّ & (71) تَخَلّٰى & (89) اِجْتَبٰى & \\
(3) اِسْلَنْقٰى & (3) أَدَبَ & (30) خَافَ & (54) أَوْجَزَ & (72) تَوَفّٰى & (90) اِبْتَلٰى & \\
\thead{باب نَصَرَ} & (10) هَنَأَ & (33) جَيِدَ & (55) أَيْسَرَ & \thead{باب تَفَاعَلَ} & (91) اِتَّدٰى & \\
(1) أَخَذَ & (14) فَرَّ & (36) رَضِيَ & (56) أَعَانَ & (73) تَحَابَّ & (92) اِهْتَدٰى & \\
(13) مَدَّ & (17) وَعَدَ & (41) نَسِيَ & (57) أَطَارَ & (74) تَصَابٰى & (93) اِدَّعٰى & \\
(21) يَمَنَ & (23) يَتَمَ & (43) أَرِيَ & (58) أَسَاءَ & (75) تَسَاوٰى & \thead{باب اِسْتَفْعَلَ} & \\
 & (27) وَأَرَ & (46) وَجِيَ & (59) أَرْضٰى & \thead{باب اِنْفَعَلَ} & (94) اِسْتَأْذَنَ & \\
 & (32) بَاعَ & (49) حَيِيَ & (60) أَلْقٰى & (76) اِنْضَمَّ & (95) اِسْتَمَلَّ & \\
\end{sarf}
